% title:          Secretary problem
% area:           strategy-puzzles, discrete-probability
% methods:        symmetry
% difficulty:
% difficulty-ai:  medium
% status:
% review:         none
% --- history: all optional, leave EMPTY if unknown ---
% origin:         journal
% origin-date:    1960-02
% origin-number:
% also-in:        book
% taken-from:
% references:     Martin Gardner, Mathematical Games column 'A fifth collection of "brain-teasers"', Scientific American 202(2), February 1960, p. 150, 'The Game of Googol' (first appearance in print; game devised in 1958 by J. H. Fox Jr. and L. G. Marnie; analysis by L. Moser and J. R. Pounder) - https://www.scientificamerican.com/article/mathematical-games-1960-02/; reprinted in M. Gardner, New Mathematical Diversions from Scientific American, Simon and Schuster, 1966, Chapter 3 (cited as 'Problem 3'); reportedly introduced earlier by Merrill M. Flood as the 'fiancée problem' in a 1949 lecture (unpublished) - https://en.wikipedia.org/wiki/Secretary_problem; T. S. Ferguson, 'Who Solved the Secretary Problem?', Statistical Science 4(3), 1989, 282-289 (history incl. Cayley and Kepler) - https://projecteuclid.org/journals/statistical-science/volume-4/issue-3/Who-solved-the-secretary-problem/10.1214/ss/1177012493.full; T. S. Ferguson, Optimal Stopping and Applications, Ch. 2, §2.1 'The Classical Secretary Problem' - https://www.math.ucla.edu/~tom/Stopping/sr2.pdf; variant: Sultan's dowry problem - https://mathworld.wolfram.com/SultansDowryProblem.html; Club source: Hongler
% history:        ai
\begin{problem}
A princess has $1 \leq N$ suitors. Each suitor possesses a distinct ``score''. Among them is Prince Charming, who has the maximal score. The princess encounters the suitors in a random order. For each suitor, she must decide either to marry him or to reject him (and, if she rejects him, she cannot go back). She can compare the score of the current suitor with the scores of those she has already rejected, but otherwise she cannot recognise Prince Charming. She wishes to maximise her chance of marrying Prince Charming (if she reaches the end of the list of $N$ suitors, she must marry the last one).\\[0.5em]
\textbf{(a)} How can she find a strategy that guarantees a chance $\geq \frac{1}{4}$ of marrying Prince Charming?\\
\textbf{(b)} What is the optimal strategy? \emph{Hint:} it gives a chance $>36\%$. \emph{Hint:} the strategy is not complicated at all.\\
\end{problem}
\begin{solution}
i. Réponse: elle décide de rejeter les $N / 2$ premiers prétendants et de se marier avec le premier parmi les $N / 2$ derniers qui est mieux que les $N / 2$ premiers prétendants (s'il existe). Cette stratégie a probabilité $\geq \frac{1}{4}$ de fonctionner, car la chance que le dauphin (le préntendant qui est juste en dessous du prince charmant) soit dans le $N / 2$ premiers est $\frac{1}{2}$ et la chance que le prince charmant soit dans les $N / 2$ derniers est $\frac{1}{2}$ aussi. Comme les deux événements sont (presque) indépendants ( $N$ assez grand) la probabilité qu'ils se produisent les deux est $\approx \frac{1}{4}$ et dans ce cas de figure la princesse trouve en effet le prince charmant.\\
(b) Quelle est la stratégie optimale ? (Indice: elle donne une chance $>36 \%$ ) (Indice: la stratégie n'est pas compliquée du tout).\\
i. Stratégie: la princesse laisse tomber les $x$ premiers prétendants et se marie avec le premier qui est mieux que ces $x$ prétendants (s'il existe). En fonction de $x$, quelle est la probabilité $P(x)$ de trouver le prince charmant?\\
ii. Observation: si on regarde la suite des rangs relatifs des prétendants au moment où la princesse les rencontre, elle forme une suite dans

$$
\{1\} \times\{1,2\} \times\{1,2,3\} \times \cdots \times\{1,2, \ldots, N\}
$$

Par symétrie du problème, chaque suite de ce type est équiprobable et a une chance $\frac{1}{N!}$ de se produire.\\
iii. Dans le cas où la princesse se marie avec le prince charmant, les rangs relatifs à partir de $x$ doivent ressembler à cela: il y a un moment $p>x$ où elle voit le prince charmant (qui a donc rang relatif 1) et ensuite, si elle continuait l'expérience, elle ne verrait que des rangs relatifs $>1$.\\
iv. Quelle est la probabilité que cela se produise pour le prétendant $p$ ?

$$
\frac{x}{x+1} \frac{x+1}{x+2} \cdots \frac{p-2}{p-1} \frac{1}{p} \frac{p}{p+1} \frac{p+1}{p+2} \cdots \frac{n-1}{n}=\frac{x}{n(p-1)}
$$

v. Maintenant, il faut sommer sur tous les $p>n$ possibles. Donc on calcule

$$
\begin{aligned}
\sum_{p=x+1}^{n} \frac{x}{n(p-1)} & =\frac{x}{n} \sum_{p=x+1}^{n} \frac{1}{p-1} \\
& =\frac{x}{n} \sum_{p=x}^{n-1} \frac{1}{p} \\
& \approx \frac{x}{n}(\log (n-1)-\log (x)) \\
& =\frac{x}{n} \log \frac{n-1}{x} \approx \frac{x}{n} \log \frac{n}{x} \\
& =-\alpha \log \alpha
\end{aligned}
$$

où $x / n=\alpha$ est la fraction des prétendants que la princesse élimine d'emblée.\\
vi. Comment optimiser par rapport à $\alpha \in[0,1]$ ? Cette fonction fait 0 en 0 et en 1, est positive entre deux... Pour trouver son max, on dérive, on trouve $-\log \alpha-1$. Pour que la dérivée fasse 0 , il faut $\alpha=\frac{1}{e} \approx 37 \%$. Et la valeur de la probabilté en $\alpha=\frac{1}{e}$, est $-\frac{1}{e} \log \left(\frac{1}{e}\right)=\frac{1}{e} \approx 37 \%$.

Penser à peut etre faire une version du prophet inequality problem: \url{https://en.wikipedia.org/wiki/Prophet_inequality}
\\\\
\url{https://dl.acm.org/doi/epdf/10.1145/3219166.3219182}
\\\\
\url{https://arxiv.org/pdf/2411.01191}
\end{solution}
